\begin{lemma}[Coordinates and Gram matrix of the Hermite gradients]
Fix $n \in \mathbb{N}$ and let $D = 2^{-1/2}\nabla$ be the operator of
\noderef{GaussD}.  Then both of the following hold.

\emph{(i) (Membership and coordinate formula.)}  For every multi-index
$\alpha \in \mathbb{N}_0^n$, the Hermite function $h_\alpha$ of
\noderef{HermiteL2}, viewed as an element of $L^2(\gamma_n)$, lies in
$W^{1,2}(\gamma_n)$ in the sense of \noderef{SobolevW12}, and its energy-normalized
gradient satisfies
\[
  D h_\alpha(x)
    = \bigl(\sqrt{\alpha_1}\,h_{\alpha-e_1}(x), \ldots,
            \sqrt{\alpha_n}\,h_{\alpha-e_n}(x)\bigr)
  \qquad \text{for $\gamma_n$-almost every } x \in \mathbb{R}^n,
\]
with $h_\beta$ the Hermite function of \noderef{HermiteFun} and $\alpha - e_j$ the
multi-index obtained by decreasing the $j$th entry of $\alpha$ by one, under the
convention of \noderef{MultiIndexSum} that this entry stays at $0$ when
$\alpha_j = 0$.

\emph{(ii) (Gram matrix.)}  For all multi-indices
$\alpha, \beta \in \mathbb{N}_0^n$,
\[
  \langle D h_\alpha,\, D h_\beta\rangle_{L^2(\gamma_n;\mathbb{R}^n)}
    = \begin{cases}
        |\alpha|, & \alpha = \beta,\\
        0, & \alpha \neq \beta,
      \end{cases}
\]
where $|\alpha| = \alpha_1 + \cdots + \alpha_n$ is the order of
\noderef{MultiIndexSum}.

The equality in (i) is an identity $\gamma_n$-almost everywhere, as it must be:
$D h_\alpha$ is an element of $L^2(\gamma_n;\mathbb{R}^n)$, that is, an
equivalence class, while the right-hand side is a genuine function.  Both
assertions are stated directly about $D h_\alpha$, and not about some auxiliary
$L^2$ representative of the right-hand side; it is this that lets (ii) be used
without carrying the membership proof of (i) around with it.
\end{lemma}
\begin{proof}
Throughout, write
\[
  F_\alpha(x) = \bigl(\sqrt{\alpha_1}\,h_{\alpha-e_1}(x), \ldots,
                      \sqrt{\alpha_n}\,h_{\alpha-e_n}(x)\bigr)
  \in \mathbb{R}^n , \qquad x \in \mathbb{R}^n ,
\]
for the field appearing on the right of (i).  The proof establishes (i) by
exhibiting $\sqrt2\,F_\alpha$ as a weak Euclidean gradient of $h_\alpha$ and
appealing to \noderef{GaussDWitness}, and then reads (ii) off the coordinate
formula.

\emph{Step 1: $F_\alpha \in L^2(\gamma_n;\mathbb{R}^n)$.}
Each coordinate of $F_\alpha$ is a constant multiple of a Hermite function, which
is continuous by \noderef{HermiteFunContDiff}, so $F_\alpha$ is continuous and in
particular measurable.  For the square-integrability, bound the Euclidean norm by
the sum of the absolute values of the coordinates: for every $x$,
\[
  |F_\alpha(x)|
    = \Bigl(\sum_{j=1}^n \bigl|\sqrt{\alpha_j}\,h_{\alpha-e_j}(x)\bigr|^2\Bigr)^{1/2}
    \leq \sum_{j=1}^n \bigl|\sqrt{\alpha_j}\,h_{\alpha-e_j}(x)\bigr| ,
\]
the inequality because $\sum_j a_j^2 \leq (\sum_j a_j)^2$ for nonnegative reals
$a_j$ and the square root is monotone.  By \noderef{HermiteFunMemLp} each
$h_{\alpha-e_j}$ lies in $L^2(\gamma_n)$, hence so does each
$|\sqrt{\alpha_j}\,h_{\alpha-e_j}|$, and hence so does the finite sum on the
right.  A measurable function dominated pointwise in absolute value by an
$L^2(\gamma_n)$ function is itself in $L^2(\gamma_n)$, so
$F_\alpha \in L^2(\gamma_n;\mathbb{R}^n)$.

\emph{Step 2: $\sqrt2\,F_\alpha$ is a weak Euclidean gradient of $h_\alpha$.}
Fix $j$ and a test function $\varphi \in C_c^\infty(\mathbb{R}^n)$.  By
\noderef{HermiteFunDeriv},
\begin{equation*}
  \partial_j h_\alpha(x) = \sqrt2\,\sqrt{\alpha_j}\,h_{\alpha-e_j}(x)
    = \sqrt2\,\bigl(F_\alpha(x)\bigr)_j
  \qquad (x \in \mathbb{R}^n), \tag{1}
\end{equation*}
and $h_\alpha$ is differentiable at every point.  By
\noderef{HermiteFunContDiff}, $h_\alpha$ and $\partial_j h_\alpha$ are
continuous, and $\varphi$ and $\partial_j \varphi$ are continuous with compact
support; hence each of the three products
$(\partial_j h_\alpha)\varphi$, $h_\alpha\,\partial_j\varphi$ and
$h_\alpha\varphi$ is continuous with compact support, and therefore Lebesgue
integrable on $\mathbb{R}^n$.  With $h_\alpha$ and $\varphi$ differentiable
everywhere and these three integrability conditions in hand, the integration-by-parts
formula on $\mathbb{R}^n$ in the direction $e_j$ applies and gives
\[
  \int_{\mathbb{R}^n} h_\alpha(x)\,\partial_j\varphi(x)\,dx
    = -\int_{\mathbb{R}^n} \partial_j h_\alpha(x)\,\varphi(x)\,dx
    = -\int_{\mathbb{R}^n} \sqrt2\,\bigl(F_\alpha(x)\bigr)_j\,\varphi(x)\,dx ,
\]
the last equality by $(1)$.  As $j$ and $\varphi$ were arbitrary, this is exactly
the defining identity of \noderef{HasGaussWeakGrad} for the pair
$(h_\alpha, \sqrt2\,F_\alpha)$.  Note that these are Lebesgue integrals, as
\noderef{HasGaussWeakGrad} requires; no Gaussian weight enters at this step, and
none is needed, since all three integrands have compact support.

\emph{Step 3: proof of (i).}
By Step 1 the field $F_\alpha$ determines an element of
$L^2(\gamma_n;\mathbb{R}^n)$, which we again call $F_\alpha$; it agrees
$\gamma_n$-almost everywhere with the function $F_\alpha$, so by
\noderef{GaussWeakGradAECongr} the conclusion of Step 2 holds for that element as
well.  The same applies on the other side: $h_\alpha$ as an element of
$L^2(\gamma_n)$ via \noderef{HermiteL2} agrees $\gamma_n$-almost everywhere with
the function $h_\alpha$ of \noderef{HermiteFun}.  Thus $\sqrt2\,F_\alpha$ is a
weak Euclidean gradient of the $L^2(\gamma_n)$ element $h_\alpha$, and
\noderef{GaussDWitness} applies and yields both
\begin{equation*}
  h_\alpha \in W^{1,2}(\gamma_n)
  \qquad\text{and}\qquad
  D h_\alpha = F_\alpha \ \text{ in } L^2(\gamma_n;\mathbb{R}^n). \tag{2}
\end{equation*}
The second equality of $(2)$ is an equality of $L^2$ elements, hence says that
$Dh_\alpha$ agrees $\gamma_n$-almost everywhere with the function $F_\alpha$,
which is the coordinate formula of (i).

\emph{Step 4: proof of (ii).}
Let $\alpha,\beta \in \mathbb{N}_0^n$.  The inner product of
$L^2(\gamma_n;\mathbb{R}^n)$ pairs matching coordinates and integrates against
$\gamma_n$, so by $(2)$ applied to $\alpha$ and to $\beta$,
\begin{equation*}
  \langle Dh_\alpha, Dh_\beta\rangle
    = \sum_{j=1}^n \int_{\mathbb{R}^n}
        \bigl(F_\alpha(x)\bigr)_j \bigl(F_\beta(x)\bigr)_j \, d\gamma_n(x)
    = \sum_{j=1}^n \sqrt{\alpha_j}\sqrt{\beta_j}\;
        \langle h_{\alpha-e_j}, h_{\beta-e_j}\rangle , \tag{3}
\end{equation*}
the interchange of the finite sum with the integral being legitimate because each
summand is integrable: it is a product of two $L^2(\gamma_n)$ functions by
\noderef{HermiteFunMemLp}.  By \noderef{HermiteL2Orthonormal} the inner product in
the $j$th summand of $(3)$ equals $1$ when $\alpha - e_j = \beta - e_j$ and $0$
otherwise.

The degenerate indices must be disposed of first, because under the convention of
\noderef{MultiIndexSum} --- the $j$th entry is left at $0$ when $\alpha_j = 0$ ---
the decremented multi-indices can coincide even for $\alpha \neq \beta$: for
instance $n = 2$, $j = 1$, $\alpha = (0,5)$, $\beta = (1,5)$ give
$\alpha - e_1 = \beta - e_1 = (0,5)$.  So it is the prefactor, not the Gram entry,
that must kill those summands.

\emph{(4a) If $\alpha_j = 0$ or $\beta_j = 0$, the $j$th summand of $(3)$ is $0$.}
Its prefactor is $\sqrt{\alpha_j}\sqrt{\beta_j}$ and one of the two factors is
$\sqrt{0} = 0$, so the prefactor is $0$; the Gram entry is a real number, hence the
product vanishes whatever that entry is.

\emph{(4b) If $\alpha_j \geq 1$ and $\beta_j \geq 1$, then
$\alpha - e_j = \beta - e_j$ forces $\alpha = \beta$.}  By
\noderef{MultiIndexSum} the multi-index $\alpha - e_j$ agrees with $\alpha$ at
every entry $i \neq j$ and has $j$th entry $\alpha_j - 1$, and likewise for
$\beta$; the convention plays no role here, since both $j$th entries are at least
$1$ and the decrement is therefore genuine.  So $\alpha - e_j = \beta - e_j$ gives
$\alpha_i = \beta_i$ for every $i \neq j$ together with
$\alpha_j - 1 = \beta_j - 1$, and the latter yields $\alpha_j = \beta_j$ because
both $\alpha_j$ and $\beta_j$ are at least $1$.  Hence $\alpha = \beta$.

Suppose now $\alpha \neq \beta$, and let $j$ be arbitrary.  If $\alpha_j = 0$ or
$\beta_j = 0$, the $j$th summand is $0$ by (4a).  Otherwise $\alpha_j \geq 1$ and
$\beta_j \geq 1$, and then $\alpha - e_j \neq \beta - e_j$ --- equality would force
$\alpha = \beta$ by (4b), contradicting $\alpha \neq \beta$ --- so the Gram entry
is $0$ and the $j$th summand is $0$ again.  Every summand of $(3)$ vanishes, so
$\langle Dh_\alpha, Dh_\beta\rangle = 0$.

Suppose instead $\alpha = \beta$.  The $j$th summand of $(3)$ is
$\sqrt{\alpha_j}\sqrt{\alpha_j}\,\langle h_{\alpha-e_j}, h_{\alpha-e_j}\rangle$,
and this equals $\alpha_j$: if $\alpha_j \geq 1$ the two decremented indices are
literally the same multi-index, so the Gram entry is $1$ and
$\sqrt{\alpha_j}\sqrt{\alpha_j} = \alpha_j$; if $\alpha_j = 0$ the summand is $0$
by (4a) and $\alpha_j = 0$ as well.  Summing over $j$ and using
\noderef{MultiIndexSum} gives $\sum_{j=1}^n \alpha_j = |\alpha|$.  This is (ii).
\end{proof}
